\chapter{Pattern Recognition in Imaging}

\section*{Patient focus}
This chapter explains the examination or clinical problem in accessible language. Patients should use it to prepare questions and understand the purpose, experience, benefits, and limitations of imaging; local instructions and the official report take priority.


\label{ch:ch29_pattern_recognition}

\section{A Structured Approach}
Begin with the clinical question, compare prior studies, confirm technical adequacy, and identify the modality and sequence. Review anatomy systematically before focusing on the most conspicuous abnormality. Separate observation from interpretation and interpretation from recommendation.

\subsection{The Systematic Search Pattern}
Rather than scanning images haphazardly, adopt a consistent search pattern for every study. On a chest CT, for example, begin with the airways, then the lungs, pleura, mediastinum, heart, vessels, bones, and soft tissues. On a brain MRI, assess symmetry, then grey--white differentiation, ventricles, sulci, and extra-axial spaces. A reproducible approach reduces the chance of missing incidental or unexpected findings.

\begin{pearlbox}[Clinical Pearl: Always check the edge]
A preventable source of missed findings is incomplete image review. Lung apices, the posterior costophrenic sulcus, visible upper-abdominal structures on chest CT, and the skull vault on brain imaging can be overlooked. Build ``edge checks'' into your routine.
\end{pearlbox}

\section{Pattern Families}
Describe location, distribution, morphology, attenuation or signal, enhancement, diffusion, temporal behavior, and associated findings. Pattern recognition narrows a differential but does not replace clinical correlation or tissue diagnosis when required.

\subsection{Common Radiological Patterns}

Table~\ref{tab:pattern_families} summarises the principal pattern families encountered in everyday practice and the differential diagnoses they most commonly suggest.

\begin{table}[h]
\centering
\caption{Common imaging patterns and their principal differential diagnoses.}
\label{tab:pattern_families}
\begin{tabular}{p{3.2cm} p{4cm} p{5.5cm}}
\toprule
\textbf{Pattern} & \textbf{Distribution / Features} & \textbf{Key Differentials} \\
\midrule
Consolidation & Dense opacification obscuring vessels, air bronchograms & Pneumonia, organising pneumonia, pulmonary infarction, lymphoma \\
Cystic lung disease & Multiple thin- or thick-walled cysts & Lymphangioleiomyomatosis, pulmonary Langerhans cell histiocytosis, Birt--Hogg--Dub\'{e} syndrome \\
Diffuse bilateral infiltrates & Symmetrical or asymmetrical airspace or interstitial pattern & ARDS, pulmonary oedema, diffuse alveolar haemorrhage, miliary TB \\
Ground-glass opacity & Hazy increase in attenuation with preservation of bronchovascular margins & Pneumocystis pneumonia, viral pneumonia, pulmonary oedema, alveolar haemorrhage, adenocarcinoma in situ \\
Multiple lytic lesions & Widespread, ``punched-out'' & Multiple myeloma, metastases (thyroid, renal, breast, lung), hyperparathyroidism \\
Multiple pulmonary nodules & Random, perilymphatic, or centrilobular & Metastases, septic emboli, granulomatous disease, vasculitis \\
Solitary bone lesion & Lytic, blastic, or mixed & Metastasis, myeloma, enchondroma, fibrous dysplasia, simple bone cyst \\
Solitary pulmonary nodule & Round, $\leq 3$~cm, well or poorly defined & Granuloma, hamartoma, primary lung carcinoma, metastasis, AVM \\
\bottomrule
\end{tabular}
\end{table}

\subsection{Characterising a Lesion by Pattern}
For any lesion, systematically evaluate:
\begin{enumerate}
    \item \textbf{Location} -- organ, lobe, segment, compartment.
    \item \textbf{Morphology} -- round, irregular, spiculated, cystic, solid.
    \item \textbf{Attenuation / Signal} -- fat, water, soft tissue, air, bone, haemorrhage.
    \item \textbf{Enhancement} -- none, peripheral, nodular, homogeneous, heterogeneous.
    \item \textbf{Margins} -- well-defined, ill-defined, infiltrative.
    \item \textbf{Effect on adjacent structures} -- mass effect, bone destruction, vessel encasement.
\end{enumerate}
These attributes narrow the differential and guide the next investigation.

\begin{pearlbox}[Clinical Pearl: Interpret morphology and size together]
Nodule assessment considers size, morphology, attenuation, calcification pattern, interval growth, and patient risk. A tiny spiculated nodule and a larger calcified nodule cannot be ranked reliably from size or a generic description of calcification alone; use the full imaging and clinical context and applicable management guidance.
\end{pearlbox}

\section{Cognitive Safety}
Use a checklist, deliberately search for common missed findings, and reconsider the first impression when discordant evidence appears. Avoid anchoring, premature closure, search satisfaction, and overconfidence. Escalate uncertain or critical findings according to local policy.

\subsection{Common Cognitive Biases in Radiology}

\begin{table}[h]
\centering
\caption{Cognitive biases most relevant to imaging interpretation.}
\label{tab:cognitive_biases}
\begin{tabular}{p{3.5cm} p{9cm}}
\toprule
\textbf{Bias} & \textbf{Description and Example} \\
\midrule
Anchoring & Over-reliance on the first piece of information. A ``known'' pneumonia on the referral may cause the reader to miss an underlying mass. \\
Availability bias & Overweighting diagnoses seen recently. After reading about COVID-19 pneumonia, every ground-glass opacity is attributed to viral infection. \\
Confirmation bias & Seeking information that supports a preferred diagnosis while ignoring contradicting evidence. \\
Overconfidence & Failing to seek a second opinion or acknowledge uncertainty when clinical stakes are high. \\
Premature closure & Accepting a diagnosis before all findings are accounted for. Labeling a focal opacity as pneumonia without considering an obstructing lesion when the clinical or imaging features do not fit. \\
Search satisfaction & Finding one abnormality and stopping the search. A positive CT pulmonary angiogram may distract from an additional lung mass or pneumothorax. \\
\bottomrule
\end{tabular}
\end{table}

\subsection{Practical Steps to Reduce Error}
\begin{itemize}
    \item Always compare with prior imaging when available.
    \item Use a structured checklist for high-risk examinations (e.g.\ CT pulmonary angiography, trauma CT).
    \item Discuss uncertain cases with a colleague before finalising the report.
    \item When an unexpected finding is identified, complete the full search pattern rather than jumping to the obvious abnormality.
    \item Be explicit about uncertainty in the report rather than using vague language.
\end{itemize}

\begin{pearlbox}[Impression:]
Pattern recognition is the backbone of radiological interpretation. A systematic search pattern, awareness of common cognitive biases, and a willingness to question your first impression will substantially reduce diagnostic error.
\end{pearlbox}

\clearpage
